\documentclass[12 pt, letterpaper]{hmcpset}
\usepackage{hw-preamble}

\name{Jayden Thadani}
\class{PCMI USS 2021}
\assignment{Exercises 3}
\duedate{}

\begin{document}

\problemlist{General exercises on Witt rings}

\begin{problem}[(1)]
	What is the Witt ring of a finite field? How many isomorphism classes of anisotropic forms are there?
\end{problem}

\begin{solution} 
	There are always exactly $3$ isomorphism classes of anisotropic forms over $\mathbb{F}_{q}$ --- $\left < 1 \right >$, $\left < a \right >$ where $a$ is not a square, and $\left < 1, b \right >$ where $b$ is not in the square class of $-1$.

	Suppose $-1$ is not a square in $\mathbb{F}_{q}$. Then $W(\mathbb{F}_{q}) \cong \mathbb{Z} / (4)$ with the isomorphism given by $\left < 1 \right > \mapsto 1$, $\left < -1 \right > \to 3$, and $\left < 1, 1 \right > \mapsto 2$.

	Suppose $-1$ is a square in $\mathbb{F}_{q}$. Then $W(\mathbb{F}_{q}) \cong \mathbb{F}_{2}[x] / (1 + x^{2})$ with the isomorphism given by $\left < 1 \right > \mapsto 1$ and $\left < a \right > \mapsto x$ where $a$ is not a square. Equivalently, we have $W(\mathbb{F}_{q}) \cong \mathbb{F}_{2}[x] / (x^{2})$ with $\left < 1 \right > \mapsto 1$ and $\left < 1, a \right > \mapsto x$. 
\end{solution}

\pagebreak

\begin{problem}[(2)]
	Let $\mathbb{E} / \mathbb{F}$ be a finite extension and let $\ell : \mathbb{E} \to \mathbb{F}$ be a non-zero $\mathbb{F}$-linear map. Show that the quadratic form $V, \ell \circ q$ over $\mathbb{F}$ induces a map of Grothendieck--Witt groups (not rings) $GW(\mathbb{E}) \to GW(\mathbb{F})$ and Witt groups  $W(\mathbb{E}) \to W(\mathbb{F})$.
\end{problem}

\begin{solution} 
	Let $\ell_{*} : GW(\mathbb{E}) \to GW(\mathbb{F})$ be the map given by $(V/\mathbb{E}, q) \mapsto (V/\mathbb{F}, \ell \circ q)$ and extended linearly for additive inverses. Then
	\begin{align*}
		\ell_{*}((V, q_{V}) \oplus (W, q_{W})) & = \ell_{*}((V \oplus W)/\mathbb{E}, q_{V} + q_{W}) \\
						       & = (V \oplus W)/\mathbb{F}, \ell \circ (q_{V} + q_{W}) \\
						       & = (V/\mathbb{F}, \ell \circ q_{V}) \oplus (W/\mathbb{F}, \ell \circ q_{W}) \\
						       & = \ell_{*}(V, q_{V}) \oplus \ell_{*}(W, q_{W}).
	\end{align*}
	Thus, $\ell_{*}$ is a group homomorphism.

	Consider $H_{\mathbb{E}} = \left < 1, -1 \right >_{\mathbb{E}}$. Let $v_{1}, v_{2}$ be the $\mathbb{E}$-basis corresponding to the diagonalisation $\left < 1, -1 \right >$ and let $e_{1}, \dots, e_{n}$ be the $\mathbb{F}$-basis of $\mathbb{E} / \mathbb{F}$. Then, $\ell_{*} H_{\mathbb{E}}$ (over $\mathbb{F}$) has a diagonalisation given by the basis $v_{i} \otimes e_{j}$ with $q(v_{i} \otimes e_{j}) = q(v_{i}) \ell(e_{j})$. Each subspace $\operatorname{span}(v_{1} \otimes e_{j}, v_{2} \otimes e_{j})$ is then just $\left < \ell(e_{j}), -\ell(e_{j}) \right > \cong \left < 1, -1 \right >$. Thus, $\ell_{*} H_{\mathbb{E}}$ is a direct sum of hyperbolic planes.

	All of this is to say that the kernel of $GW(\mathbb{E}) \to W(\mathbb{E})$ is mapped into the kernel of $GW(\mathbb{F}) \to W(\mathbb{F})$. Thus, $\ell_{*} : GW(\mathbb{E}) \to GW(\mathbb{F})$ descends to a map $\ell_{*} : W(\mathbb{E}) \to W(\mathbb{F})$.
\end{solution}

\pagebreak

\begin{problem}[(3)]
	If $\mathbb{E} / \mathbb{F}$ is any field extension, then check that there is a homomorphism of Witt rings $W(\mathbb{F}) \to W(\mathbb{E})$ given by extension of scalars of quadratic forms. If $\mathbb{E} = \mathbb{F}(\sqrt{d})$ then show that the kernel of $W(\mathbb{F}) \to W(\mathbb{E})$ is the ideal generated by $\left < 1, -d \right >$.
\end{problem}

\begin{solution}
	By induction, assume that for every $(U, q) \in \ker \varphi$ with $\dim U < n$, is in the ideal generated by $\left < 1, -d \right >$.

	Suppose $V, q$ is an $n$-dimensional anisotropic quadratic form over $\mathbb{F}$, in the kernel of $W(\mathbb{F}) \to W(\mathbb{E})$. Then $V_{\mathbb{E}}, q_{\mathbb{E}} = H_{\mathbb{E}}^{\oplus n}$. In particular, it is isotropic, and by exercises 2, problem (2), we must have $V, q \cong (U, q_{\mid U}) \oplus \left < a, -ad \right >$. Since $\left < a, -ad \right > \cong \left < a \right > \otimes \left < 1, -d \right >$, the right summand is in the ideal generated by $\left < 1, -d \right >$. Since $\left < a, -ad \right > \in \ker \varphi$ we also have $U, q_{\mid U} \in \ker \varphi$. Then, by induction, $U, q_{\mid U}$ is also in the ideal generated by $\left < 1, -d \right >$.
\end{solution}

\pagebreak

\begin{problem}[(4)]
	Let $\mathbb{E} / \mathbb{F}$ be a cubic extension. We will show that a quadratic form over $\mathbb{F}$, $V, q \cong \left < a_{1}, \dots, a_{n} \right >$ with isotropic base change to $\mathbb{E}$ must be isotropic itself. 

	Let $\alpha \in \mathbb{E}$ be a primitive element with minimal polynomial $p(x) \in \mathbb{F}[x]$. Show that we can find polynomials  $f_{i}(x) \in \mathbb{F}[x]$ of degree at most $2$ such that $p(x) \mid \sum_{i = 1}^{n} a_{i} f_{i}(x)^{2}$. Show that the right hand side necessarily has a root in $\mathbb{F}$ and that $V, q$ is therefore isotropic over $\mathbb{F}$. Deduce that the map $W(\mathbb{F}) \to W(\mathbb{E})$ is injective in this case.
\end{problem}

\begin{solution}
	Suppose $\mathbb{E} = \mathbb{F}(\alpha)$ and $V_{\mathbb{E}}$ is isotropic. Consider non-zero $v' \in V_{\mathbb{E}}$ with $q_{\mathbb{E}}(v) = 0$. Then we can write each $v_{i}' = f_{i}(\alpha)$ for some polynomials $f_{i} \in \mathbb{F}[x]$ of degree at most $2$ and get $g(x) = \sum_{i = 1}^{n} a_{i} f_{i}(\alpha)^{2} = 0$.

	Then, by its definition, $p(x)$ divides $g(x)$. Since $p$ is degree $3$, $g$ must be degree $4$ (it is of even degree at most $4$ and at least $3$) and their quotient is degree $1$. Thus,  $g(x) = p(x)(x - \lambda)$ for some $\lambda \in \mathbb{F}$. That is, $g$ has a root over $\mathbb{F}$, and so $V, q$ is isotropic over $\mathbb{F}$.

	Let $\varphi : W(\mathbb{F}) \to W(\mathbb{E})$ be the base change homomorphism. We show $\ker \varphi = 0$. Suppose $\varphi : V, q \mapsto 0$ and $V, q \cong (U, q_{\mid U}) \oplus H^{\oplus r}$ is the decomposition into a hyperbolic part and a (possibly zero) anisotropic part. Then $U, q_{\mid U} \mapsto 0$. That is, the base change to $\mathbb{E}$ is zero or hyperbolic and isotropic. $U, q$ is not isotropic over $\mathbb{F}$, so its base change to $\mathbb{E}$ is not isotropic. Thus, $U, q_{\mid U} = 0$. That is, $V, q$ is hyperbolic and thus, $0$ in the $W(\mathbb{F})$.
\end{solution}

\pagebreak

\problemlist{Exercises on the Witt ring and orderings}

\begin{problem}[(1)]
	Let $d \in \mathbb{F}^{\times}$. Show that the kernel of $W(\mathbb{F}) \to W(\mathbb{F}(\sqrt{d}))$ is the ideal of $W(\mathbb{F})$ generated by $\left < 1, -d \right >$.
\end{problem}

\pagebreak

\begin{problem}[(2)]
	All torsion in the Witt ring $W(\mathbb{F})$ is $2$-power torsion (due to Pfister, Scharlau). Prove this as follows:
	\begin{enumerate}
		\item Let $\alpha, \beta \in \mathbb{F}^{\times}$ such that $\alpha^{2} + \beta^{2} \neq 0$. Then $\left < 1, -(\alpha^{2} + \beta^{2}) \right >$ is $2$-torsion.

		\item Let $\alpha, \beta \in \mathbb{F}^{\times}$ such that $\alpha^{2} + \beta^{2} \neq 0$. Then the kernel of $W(\mathbb{F}) \to W(\mathbb{F}(\sqrt{\alpha^{2} + \beta^{2}}))$ consists of $2$-torsion.

		\item Let $\mathbb{F}$ be a field such that any sum of squares is a square ($\mathbb{F}$ is a \emph{pythagorean field}). Then either $\mathbb{F}$ is quadratically closed and $W(\mathbb{F}) \cong \mathbb{Z} / (2)$ or $W(\mathbb{F})$ is torsion free.

		\item Combine the previous three parts to deduce the claim.
	\end{enumerate}
\end{problem}

\begin{solution}
	Let $d = \alpha^{2} + \beta^{2}$.
	\begin{enumerate}
		\item By exercises 1, problem (4) we have
			\[
				\left < 1, -d \right > \oplus \left < 1, -d \right > = H^{\oplus 2} = 0 \in W(\mathbb{F}).
			\]

		\item By the previous exercise, the kernel of  $W(\mathbb{F}) \to W(\mathbb{F}(\sqrt{d}))$ is the ideal generated by $\left < 1, -d \right >$. Since (for our choice of $d$) $\left < 1, -d \right >$ is $2$-torsion, the result follows.

		\item Suppose $\mathbb{F}$ is pythagorean and suppose $(V, q) = \left < a_{1}, \dots, a_{n} \right >$ is an $n$-dimensional form with $(V, q)^{\oplus m}$ isotropic. Then there exist some $x_{i, j}$ for each  $1 \le i \le m$ and $1 \le j \le n$ such that
			\[
				\sum_{i = 1}^{m} q(x_{i}) = \sum_{j = 1}^{n} a_{j} \left ( \sum_{i = 1}^{m} x_{i, j}^{2} \right ) = \sum_{j = 1}^{n} a_{j} y_{j}^{2} = q(y)
			\]
			Here we write $x_{i} = (x_{i, 1}, \dots, x_{i, n}) \in V$ in the basis corresponding to the diagonalisation and $y = (y_{1}, \dots, y_{n}) \in \mathbb{F}^{\oplus n}$ with $y_{j}^{2} = \sum_{i = 1}^{m} x_{i, j}^{2}$ for each $j$, by the pythagorean property of $\mathbb{F}$. Then either $y \neq 0$ and $V, q$ is isotropic or $y = 0$ and $\left < 1 \right >^{\oplus m}$ is isotropic.

			If $\left < 1 \right >^{\oplus m}$ is isotropic for any $m \ge 2$, then $-1$ is a sum of squares. Then since every element is a difference of squares, every element becomes a sum of squares, and thus, a square. That is, if $\left < 1 \right >^{\oplus m}$ is isotropic, $\mathbb{F}$ is quadratically closed. In this case we know $W(\mathbb{F}) \cong \mathbb{Z} / (2)$ with isomorphism given by ``dimension modulo $2$''.

			Else, there are no anisotropic forms $V, q$ with isotropic direct sums $(V, q)^{\oplus m}$. Thus, there is no torsion $W(\mathbb{F})$.

		\item It is a general fact that every field $\mathbb{F}$ has a pythagorean closure $\mathbb{F}^{\text{pyth}}$ (the smallest pythagorean field containing it). Fixing an algebraic closure $\mathbb{F} \subseteq \overline{\mathbb{F}}$, we can consider $\mathbb{F}^{\text{pyth}}$ to be the union of all algebraic extensions $\mathbb{K} \subseteq \overline{\mathbb{F}}$ that are themselves towers $\mathbb{F} \subseteq \mathbb{K}_{1} \subseteq \dots \subseteq \mathbb{K}_{n} = \mathbb{K}$ with $\mathbb{K}_{i + 1} = \mathbb{K}_{i}(\sqrt{\alpha_{i}^{2} + \beta_{i}^{2}})$.

 Let $\mathbb{F}$ be an arbitrary field of characteristic not $2$. If $\mathbb{F}$ is pythagorean, we are done. Else, consider the map $W(\mathbb{F}) \to W(\mathbb{F}^{\text{pyth}})$. Since $W(\mathbb{F}^{\text{pyth}})$ is either $\mathbb{Z} / (2)$ or torsion free, all torsion lies in the kernel of this map or is $2$-power torsion. Explicitly, if $n x = 0$ for $n = p_{1}^{m_{1}} \cdots p_{k}^{m_{k}}$, then either one of the $p_{i} = 2$ or one of the $p_{i}$ is in the kernel. Choosing minimal $n$ gives the desired result.

			By a general trick about towers of field extensions and finitary conditions, any element of the kernel of the map $W(\mathbb{F}) \to W(\mathbb{F}^{\text{pyth}})$ must lie in the kernel of some $W(\mathbb{F}) \to W(\mathbb{K})$. Thus, all torsion in $W(\mathbb{F})$ is $2$-power torsion.
	\end{enumerate}
\end{solution}

\pagebreak

\begin{problem}[(3)]
	Show that the Witt ring $W(\mathbb{F})$ is torsion if and only if $-1$ is a sum of squares in $\mathbb{F}$.
\end{problem}

\begin{solution}
	If $W(\mathbb{F})$ is torsion, then $W(\mathbb{F}^{\text{pyth}})$ is torsion. Thus, $\mathbb{F}^{\text{pyth}}$ is quadratically closed. In particular $-1 \in \mathbb{F}^{\text{pyth}}$ is a square and so $-1$ is a sum of squares in $\mathbb{F}$.
\end{solution}

\pagebreak

\begin{problem}[(4)]
	The signature map gives an isomorphism $W(\mathbb{R}) \cong \mathbb{Z}$. More generally, given any ordered field $\mathbb{F}$, show that we obtain a ring homomorphism $\sigma : W(\mathbb{F}) \to \mathbb{Z}$ given by taking the \emph{generalised signature} of a quadratic form.

	Conversely, show that each ring homomorphism $\varphi : W(\mathbb{F}) \to \mathbb{Z}$ gives an ordering on $\mathbb{F}$ as follows.

	\begin{enumerate}
		\item Suppose $\varphi : W(\mathbb{F}) \to \mathbb{Z}$ is a ring homomorphism. Define a subset $P \subseteq \mathbb{F}^{\times}$ of \emph{positive elements} as follows;  $\alpha \in \mathbb{F}^{\times}$ belongs to $P$ if $\varphi(\left < \alpha \right >) = 1$. Note in fact that for any $\alpha \in \mathbb{F}^{\times}$ we have $\varphi(\left < \alpha \right >) = \pm 1$.

		\item Show that $P$ is closed under sums and products. Moreover, show that $\mathbb{F}^{\times} = P \sqcup (-P)$ is a disjoint union.
	\end{enumerate}
\end{problem}

\begin{solution}
	If $<$ gives an ordering on $\mathbb{F}$, we claim it defines a homomorphism $\sigma : W(\mathbb{F}) \to \mathbb{Z}$ generated by  $\sigma(\left < \alpha \right >) = 1$ if $\alpha > 0$ and $\sigma(\left < \alpha \right >) = -1$ if $\alpha < 0$. 

	To show that this is well defined, it suffices to check that multiplication and all of Witt's chain equivalences are preserved by the homomorphism. $\sigma(\left < \alpha \beta \right >) = \sigma(\left < \alpha \right >) \sigma(\left < \beta \right >)$ by the definition of an ordered field. It follows that $\sigma(\left < \alpha u^{2} \right >) = \sigma(\left < \alpha \right >)$ for $\alpha, u \in \mathbb{F}^{\times}$ since $(\pm 1)^{2} = 1$. Clearly, $\sigma(\left < \alpha, \beta \right >) = \sigma(\left < \beta, \alpha \right >)$ since $\mathbb{Z}$ is commutative. Finally we need to show
	 \[
		\sigma \left ( \left < \alpha, \beta \right > \right ) = \sigma  \left ( \left < \alpha + \beta, \frac{\alpha \beta}{\alpha + \beta} \right > \right ).
	\]
	If both $\alpha$ and $\beta$ are positive or both are negative, this follows easily. Without loss of generality, assume $\alpha > 0, \beta < 0$ and note $\sigma(\left < \alpha, \beta \right >) = 0$. If $\alpha + \beta > 0$, then the right side of the equation is $1 + (-1) / 1 = 0$. If $\alpha + \beta < 0 $ then the right hand side is $(-1) + (-1) / (-1) = 0$.

	\begin{enumerate}
		\item For $\alpha \in \mathbb{F}^{\times}$ we have $\left < \alpha \right > \otimes \left < 1 / \alpha \right > = \left < 1 \right >$. Since $\left < 1 \right > \mapsto 1$ under any ring homomorphism $W(\mathbb{F}) \to \mathbb{Z}$, we must have that $\left < \alpha \right >$ maps to a unit.

		\item Suppose $\alpha, \beta \in P$. Then $\varphi(\left < \alpha \beta \right >) = \varphi( \left < \alpha \right > \otimes \left < \beta \right >) = 1 \times 1$ and so $\alpha \beta \in P$. Also, since $\left < \alpha, \beta \right > = \left < \alpha + \beta, \alpha \beta / (\alpha + \beta) \right >$
			\[
				\varphi(\left < \alpha + \beta \right >) = \varphi(\left < \alpha \right >) + \varphi(\left < \beta \right >) - \varphi(\left < \alpha \beta / (\alpha + \beta) \right >) = 1 + 1 - (\pm 1).
			\]
			Since $\varphi(\left < \alpha + \beta \right >) = \pm 1$, it must be $1$ and $\alpha + \beta \in P$.

			If $\alpha \in P$, then $\left < \alpha \right > \oplus \left < -\alpha \right > = 0 \in W(\mathbb{F})$ so $\varphi(\left < -\alpha \right >) = -1$ and $-\alpha \notin P$. That is, $P \cap (-P) = \text{\O}$.
	\end{enumerate}

	Thus, we can define an ordering on $\mathbb{F}$ by  $\alpha > \beta$ if $\alpha - \beta \in P$. If $\alpha > \beta$, then clearly $\alpha + \gamma > \beta + \gamma$. Also, if $\alpha, \beta > 0$, then  $\alpha \beta > 0$.
\end{solution}

\pagebreak

\begin{problem}[(5) Pfister's local-to-global principle]
	Given an element $x \in W(\mathbb{F})$, $x$ is torsion if and only if its signature relative to each ordering of $\mathbb{F}$ vanishes, or equivalently, if $x$ is annihilated by all homomorphisms $W(\mathbb{F}) \to \mathbb{Z}$. Prove this as follows:
	\begin{enumerate}
		\item Suppose $x$ is non-torsion. Using Zorn's lemma show that, without loss of generality, we can assume $x$ maps to a torsion element in $W(\mathbb{F}')$ for all proper extensions $\mathbb{F}' / \mathbb{F}$.

		\item Show that for each $d \in \mathbb{F}^{\times} \setminus \mathbb{F}^{\times 2}$, $2^{n} x$ is divisible by $\left < 1, -d \right >$ in $W(\mathbb{F})$ for $n \gg 0$. Conclude that $2^{n} x = 2^{n} x \left < -d \right >$ in $W(\mathbb{F})$.

		\item Suppose that an element $d \in \mathbb{F}^{\times}$ has the property that neither $d$ nor $-d$ is a square. Then show that $2^{n} x \left < d \right > = 2^{n} x \left < -d \right >$ in $W(\mathbb{F})$. Since, for any $a \in \mathbb{F}^{\times}$, we have $\left < a \right > = - \left < - a \right >$ in $W(\mathbb{F})$, we conclude that $2^{n + 1} x = 0$ in $W(\mathbb{F})$.

			Conclude that $\mathbb{F}^{\times}$ is the (disjoint) union of the squares and the negative of the squares.

		\item Conclude that the squares give an ordering on $\mathbb{F}$ and in fact $W(\mathbb{F}) = \mathbb{Z}$, completing the proof of Pfister's local-to-global theorem.

		\item As an example, we recover the following classical fact:  $\mathbb{F}$ admits an ordering if and only if $-1$ is not a sum of squares in $\mathbb{F}$.
	\end{enumerate}
\end{problem}

\begin{solution}
	It's clear that if $x \in W(\mathbb{F})$ is torsion, then $x \mapsto 0$ under every ring homomorphism $W(\mathbb{F}) \to \mathbb{Z}$. Below, we prove (the contrapositive of) the converse. Suppose $x \in W(\mathbb{F})$ is non-torsion.

	\begin{enumerate} 
		\item Consider the poset of all field extensions $\mathbb{K} / \mathbb{F}$ such that $x_{\mathbb{K}}$ is non-torsion in $W(\mathbb{K})$. 

			Let $\{ \mathbb{K}_{\alpha} \}_{\alpha \in i}$ be a chain of such fields. Consider $\mathbb{K} = \bigcup_{\alpha \in i} \mathbb{K}_{\alpha}$ contains all of the $\mathbb{K}_{\alpha}$. By a general trick about towers of field extensions and finitary conditions, if $x_{\mathbb{K}}$ were torsion, then there must be some $\alpha$ such that $x_{\mathbb{K}_{\alpha}}$ is torsion.

			Thus, by Zorn's lemma there is a maximal $\mathbb{L} / \mathbb{F}$ such that $x_{\mathbb{L}}$ is non-torsion. It suffices to show that there is a homomorphism $W(\mathbb{L}) \to \mathbb{Z}$ that doesn't annihilate $x_{\mathbb{L}}$ --- composition with the base change $W(\mathbb{F}) \to W(\mathbb{L})$ gives the desired homomorphism.

			From here on, assume $\mathbb{F} = \mathbb{L}$.

		\item Let $\varphi : W(\mathbb{F}) \to W(\mathbb{F}(\sqrt{d}))$ be the base change homomorphism. $\ker \varphi$ is generated by $\left < 1, -d \right >$. By assumption, $x$ is torsion in $W(\mathbb{F}(\sqrt{d}))$, and necessarily $2$-power torsion. Thus, $2^{n} x = \left < 1 -d \right > y \in \ker \varphi$ for sufficiently large $n$.

		\item Suppose by way of contradiction $\pm d$ are both non-square units. Then $2^{n} x\left < d \right > = 2^{n} x = 2^{n} x \left < -d \right >$. It follows that $2^{n + 1} x \left < d \right > = 0$. $\left < d \right >$ is a unit, so $x$ must be $2^{n + 1}$-torsion in $W(\mathbb{F})$. This contradicts our hypothesis that $x$ is non-torsion in $\mathbb{F}$. Thus, one of $\pm d$ is a square and $\mathbb{F}^{\times}$ is the union of squares and non-squares.

			This union is disjoint. Else $-1$ would be a square, so $\left < 1 \right >$ would be torsion, so $W(\mathbb{F})$ would be torsion.

		\item Let $\sigma : W(\mathbb{F}) \to \mathbb{Z}$ be the homomorphism corresponding to the ordering on $\mathbb{F}$ by squares. $\sigma$ is surjective --- we have $\left < 1 \right >^{\oplus n} \mapsto n$ for each $n \in \mathbb{Z}$ (writing $\left < 1 \right >^{-1} = \left < -1 \right >$). 

		\item If $-1$ is a sum of squares in $\mathbb{F}$, then $W(\mathbb{F})$ is torsion and admits no orderings. If $-1$ is not a sum of squares, then $\left < 1 \right >$ is not torsion and so $W(\mathbb{F})$ admits a homomorphism to $\mathbb{Z}$ that doesn't annihilate $\left < 1 \right >$.
	\end{enumerate}
\end{solution}

\pagebreak

\begin{problem}[(6)]
	Let $\gamma \in \mathbb{F}^{\times}$. Show that $\gamma$ is a sum of squares if and only if $\left < 1, -\gamma \right >$ is a torsion element of the Witt ring if and only if $\gamma$ is positive with respect to any ordering of $\mathbb{F}$.
\end{problem}

\begin{solution}
	Suppose $\gamma$ is positive with respect to each ordering. Then each homomorphism $\varphi : W(\mathbb{F}) \to \mathbb{Z}$ has $\left < \gamma \right > \mapsto 1$. Thus, $\left < 1, - \gamma \right > = \left < 1 \right > - \left < \gamma \right > \in \ker \varphi$. It follows that $\left < 1, -\gamma \right >$ is torsion.

	Suppose $\left < 1, - \gamma \right >$ is torsion in $W(\mathbb{F})$. Then there is some $n$ such that $\left < 1, -\gamma \right >^{\oplus n} \cong H^{\oplus n}$. Adding a direct summand of $\left < \gamma \right >^{\oplus n - 1}$ on both sides we get $\left < 1 \right >^{\oplus n} \oplus \left < -\gamma \right > \oplus H^{\oplus n - 1} \cong H^{\oplus n} \oplus \left < \gamma \right >^{\oplus n - 1}$. Finally, since the monoid of quadratic forms is cancellative, we get $\left < 1 \right >^{\oplus n} \oplus \left < - \gamma \right > \cong H \oplus \left < \gamma \right >^{\oplus n - 1}$. Thus, $\left < 1 \right >^{\oplus n} \oplus \left < -\gamma \right >$ is isotropic, and $\gamma$ is a sum of squares.

	Suppose $\gamma$ is a sum of squares. Squares are positive in every ordering and so is $\gamma$.
\end{solution}

\end{document}
